%Root main.tex
%---------------------------------------------------------------------------
%付録A　付録
%---------------------------------------------------------------------------

\lstnewenvironment{mylisting}[1][]
    {\lstset{
        frame=single,
        numbersep=10pt,
        tabsize=2,
        #1
    }
}{}
\begin{mylisting}[language=c,caption=SRDCDB制御モジュール,label=hurokuB_SRDBcontrol]

#include <Stdlib.h>
#include <String.h>
#include <math.h>
#include <Psim.h>

// PLACE GLOBAL VARIABLES OR USER FUNCTIONS HERE...
double wp,zp,a,b;
double z11,z12,z13,g1;
double E,e,T;
double R,L,C;
double A0;
double ILa,ILb;
double Va,Vb;
double IdisOa,IdisOb;
double ua,ub;
double ILu,ILv;
double Vuv,Vvw;
double IOu,IOv;
double IdisOu,IdisOv;
double uu,uv,uw;
double Vuvref,Vvwref,Varef,Vbref;

double fc,fs;

/////////////////////////////////////////////////////////////////////
// FUNCTION: SimulationStep
//   This function runs at every time step.
//double t: (read only) time
//double delt: (read only) time step as in Simulation control
//double *in: (read only) zero based array of input values. in[0] is the first node, in[1] second input...
//double *out: (write only) zero based array of output values. out[0] is the first node, out[1] second output...
//int *pnError: (write only)  assign  *pnError = 1;  if there is an error and set the error message in szErrorMsg
//    strcpy(szErrorMsg, "Error message here..."); 
// DO NOT CHANGE THE NAME OR PARAMETERS OF THIS FUNCTION
void SimulationStep(
		double t, double delt, double *in, double *out,
		 int *pnError, char * szErrorMsg,
		 void ** reserved_UserData, int reserved_ThreadIndex, void * reserved_AppPtr)
{
// ENTER YOUR CODE HERE...
ILu = in[0];
ILv =  in[1];
IOu = in[2];
IOv = in[3];
Vuv = in[4];
Vvw = in[5];
Vuvref = in[6];
Vvwref = in[7];
e = in[8];
R = in[9];
L = in[10];
C = in[11];
T= in[12];

/****外乱電流の導出****/
IdisOu = IOu - ((2.0 * Vuv + Vvw)/R);
IdisOv = IOv + ((Vuv - Vvw)/R);


/****三相二相変換****/
ILa = ILu;
ILb = (1.0 / sqrt(3.0)) * ILu + (2.0 / sqrt(3.0)) * ILv;
IdisOa = IdisOu;
IdisOb = (1.0 / sqrt(3.0)) * IdisOu + (2.0 / sqrt(3.0)) * IdisOv;
Va = (2.0 / 3.0) * Vuv + (1.0 / 3.0) * Vvw;
Vb = (1.0 / sqrt(3.0)) * Vvw; 
Varef = (2.0 / 3.0) * Vuvref + (1.0 / 3.0) * Vvwref;
Vbref = (1.0 / sqrt(3.0)) * Vvwref; 


/****デッドビートゲインの導出****/
//E = 0.5 * e;
E = e;
wp = 1.0 / (sqrt(L*C));
zp = L / (2.0 * R * sqrt(L*C));
a = (-1.0) * wp * zp;
b = wp * sqrt((1.0 / 3.0) - (zp * zp));

z11 = exp(a * T) * (cos(b * T) + (a / b) * sin(b * T));
z12 = (1.0 / (3.0 * C * b)) * exp(a * T) * sin(b * T);
z13 = (-1.0) * (1.0 / (3.0 * C * b)) * exp(a * T) * sin(b * T);
g1 = (E / (3.0 * L * C * b)) * exp((a * T) * 0.5) * sin((b * T) * 0.5);
A0  = 1.0 /g1;


/****操作量の導出****/
ua = (A0 * Varef) - (A0 * z11 * Va) - (A0 * z12 * ILa) - (A0 * z13 * IdisOa);
ub = (A0 * Vbref) - (A0 * z11 * Vb) - (A0 * z12 * ILb) - (A0 * z13 * IdisOb);


/*output*/
out[0] = ua;
out[1] = ub;

/*read only*/
out[2] = IdisOu;
out[3] = IdisOv;
out[4] = wp;
out[5] = zp;
out[6] = a;
out[7] = b;
out[8] = Varef;
out[9] = Va;
out[10] = ILa;
out[11] = IdisOa;
out[12] = wp;
out[13] = zp;
out[14] = a;
out[15] = b;
out[16] = z11;
out[17] = z12;
out[18] = z13;
out[19] = g1;

}


/////////////////////////////////////////////////////////////////////
// FUNCTION: SimulationBegin
//   Initialization function. This function runs once at the beginning of simulation
//   For parameter sweep or AC sweep simulation, this function runs at the beginning of each simulation cycle.
//   Use this function to initialize static or global variables.
//const char *szId: (read only) Name of the C-block 
//int nInputCount: (read only) Number of input nodes
//int nOutputCount: (read only) Number of output nodes
//int nParameterCount: (read only) Number of parameters is always zero for C-Blocks.  Ignore nParameterCount and pszParameters
//int *pnError: (write only)  assign  *pnError = 1;  if there is an error and set the error message in szErrorMsg
//    strcpy(szErrorMsg, "Error message here..."); 
// DO NOT CHANGE THE NAME OR PARAMETERS OF THIS FUNCTION
void SimulationBegin(
		const char *szId, int nInputCount, int nOutputCount,
		 int nParameterCount, const char ** pszParameters,
		 int *pnError, char * szErrorMsg,
		 void ** reserved_UserData, int reserved_ThreadIndex, void * reserved_AppPtr)
{
// ENTER INITIALIZATION CODE HERE...




}


/////////////////////////////////////////////////////////////////////
// FUNCTION: SimulationEnd
//   Termination function. This function runs once at the end of simulation
//   For parameter sweep or AC sweep simulation, this function runs at the end of each simulation cycle.
//   Use this function to de-allocate any allocated memory or to save the result of simulation in an alternate file.
// Ignore all parameters for C-block 
// DO NOT CHANGE THE NAME OR PARAMETERS OF THIS FUNCTION
void SimulationEnd(const char *szId, void ** reserved_UserData, int reserved_ThreadIndex, void * reserved_AppPtr)
{




}



\end{mylisting}

\begin{mylisting}[language=c,caption=三相二相変換,label=hurokuB_32]

#include <Stdlib.h>
#include <String.h>
#include <math.h>
#include <Psim.h>

// PLACE GLOBAL VARIABLES OR USER FUNCTIONS HERE...

double a,b,c;
double aa,bb;;

//
double s2,s3,s4;

/////////////////////////////////////////////////////////////////////
// FUNCTION: SimulationStep
//   This function runs at every time step.
//double t: (read only) time
//double delt: (read only) time step as in Simulation control
//double *in: (read only) zero based array of input values. in[0] is the first node, in[1] second input...
//double *out: (write only) zero based array of output values. out[0] is the first node, out[1] second output...
//int *pnError: (write only)  assign  *pnError = 1;  if there is an error and set the error message in szErrorMsg
//    strcpy(szErrorMsg, "Error message here..."); 
// DO NOT CHANGE THE NAME OR PARAMETERS OF THIS FUNCTION
void SimulationStep(
		double t, double delt, double *in, double *out,
		 int *pnError, char * szErrorMsg,
		 void ** reserved_UserData, int reserved_ThreadIndex, void * reserved_AppPtr)
{
// ENTER YOUR CODE HERE...
//in
a = in[0];
b = in[1];
c = in[2];

//3相線間2相相対変換
//keisan
s2 = sqrt(3.0)/2.0;
s3 = 1.0/2.0;
s4 = 2.0/3.0;

//aa = s4*(s3*(a-c));
//bb = s4*(s2*b);

aa = (a/3.0)-(c/3.0);
bb = b/sqrt(3.0);


//aa = s3*(a -s2*b -s2*c);
//bb = s3*(sqrt(3.0)/2)*(b-c);


//out
out[0] = aa;
out[1] = bb;



}


/////////////////////////////////////////////////////////////////////
// FUNCTION: SimulationBegin
//   Initialization function. This function runs once at the beginning of simulation
//   For parameter sweep or AC sweep simulation, this function runs at the beginning of each simulation cycle.
//   Use this function to initialize static or global variables.
//const char *szId: (read only) Name of the C-block 
//int nInputCount: (read only) Number of input nodes
//int nOutputCount: (read only) Number of output nodes
//int nParameterCount: (read only) Number of parameters is always zero for C-Blocks.  Ignore nParameterCount and pszParameters
//int *pnError: (write only)  assign  *pnError = 1;  if there is an error and set the error message in szErrorMsg
//    strcpy(szErrorMsg, "Error message here..."); 
// DO NOT CHANGE THE NAME OR PARAMETERS OF THIS FUNCTION
void SimulationBegin(
		const char *szId, int nInputCount, int nOutputCount,
		 int nParameterCount, const char ** pszParameters,
		 int *pnError, char * szErrorMsg,
		 void ** reserved_UserData, int reserved_ThreadIndex, void * reserved_AppPtr)
{
// ENTER INITIALIZATION CODE HERE...




}


/////////////////////////////////////////////////////////////////////
// FUNCTION: SimulationEnd
//   Termination function. This function runs once at the end of simulation
//   For parameter sweep or AC sweep simulation, this function runs at the end of each simulation cycle.
//   Use this function to de-allocate any allocated memory or to save the result of simulation in an alternate file.
// Ignore all parameters for C-block 
// DO NOT CHANGE THE NAME OR PARAMETERS OF THIS FUNCTION
void SimulationEnd(const char *szId, void ** reserved_UserData, int reserved_ThreadIndex, void * reserved_AppPtr)
{




}



\end{mylisting}

\begin{mylisting}[language=c,caption=三相線間二相変換,label=hurokuB_3_2]

#include <Stdlib.h>
#include <String.h>
#include <math.h>
#include <Psim.h>

// PLACE GLOBAL VARIABLES OR USER FUNCTIONS HERE...

double a,b,c;
double aa,bb;;

//
double s2,s3,s4;

/////////////////////////////////////////////////////////////////////
// FUNCTION: SimulationStep
//   This function runs at every time step.
//double t: (read only) time
//double delt: (read only) time step as in Simulation control
//double *in: (read only) zero based array of input values. in[0] is the first node, in[1] second input...
//double *out: (write only) zero based array of output values. out[0] is the first node, out[1] second output...
//int *pnError: (write only)  assign  *pnError = 1;  if there is an error and set the error message in szErrorMsg
//    strcpy(szErrorMsg, "Error message here..."); 
// DO NOT CHANGE THE NAME OR PARAMETERS OF THIS FUNCTION
void SimulationStep(
		double t, double delt, double *in, double *out,
		 int *pnError, char * szErrorMsg,
		 void ** reserved_UserData, int reserved_ThreadIndex, void * reserved_AppPtr)
{
// ENTER YOUR CODE HERE...
//in
a = in[0];
b = in[1];
c = in[2];

//3相線間2相相対変換
//keisan
s2 = sqrt(3.0)/2.0;
s3 = 1.0/2.0;
s4 = 2.0/3.0;

//aa = s4*(s3*(a-c));
//bb = s4*(s2*b);

aa = (a/3.0)-(c/3.0);
bb = b/sqrt(3.0);


//aa = s3*(a -s2*b -s2*c);
//bb = s3*(sqrt(3.0)/2)*(b-c);


//out
out[0] = aa;
out[1] = bb;



}


/////////////////////////////////////////////////////////////////////
// FUNCTION: SimulationBegin
//   Initialization function. This function runs once at the beginning of simulation
//   For parameter sweep or AC sweep simulation, this function runs at the beginning of each simulation cycle.
//   Use this function to initialize static or global variables.
//const char *szId: (read only) Name of the C-block 
//int nInputCount: (read only) Number of input nodes
//int nOutputCount: (read only) Number of output nodes
//int nParameterCount: (read only) Number of parameters is always zero for C-Blocks.  Ignore nParameterCount and pszParameters
//int *pnError: (write only)  assign  *pnError = 1;  if there is an error and set the error message in szErrorMsg
//    strcpy(szErrorMsg, "Error message here..."); 
// DO NOT CHANGE THE NAME OR PARAMETERS OF THIS FUNCTION
void SimulationBegin(
		const char *szId, int nInputCount, int nOutputCount,
		 int nParameterCount, const char ** pszParameters,
		 int *pnError, char * szErrorMsg,
		 void ** reserved_UserData, int reserved_ThreadIndex, void * reserved_AppPtr)
{
// ENTER INITIALIZATION CODE HERE...




}


/////////////////////////////////////////////////////////////////////
// FUNCTION: SimulationEnd
//   Termination function. This function runs once at the end of simulation
//   For parameter sweep or AC sweep simulation, this function runs at the end of each simulation cycle.
//   Use this function to de-allocate any allocated memory or to save the result of simulation in an alternate file.
// Ignore all parameters for C-block 
// DO NOT CHANGE THE NAME OR PARAMETERS OF THIS FUNCTION
void SimulationEnd(const char *szId, void ** reserved_UserData, int reserved_ThreadIndex, void * reserved_AppPtr)
{




}



\end{mylisting}